\scanpage{174}

Il en résulte en effet que l'on a une suite exacte (en se bornant aux faisceaux inversibles)

\[
0\longrightarrow\mathrm{Pic}(S)\longrightarrow\mathrm{Pic}(X)\longrightarrow\mathrm{Prépic}'(X/S)
\]

% slanted
\margin{[surjectif si $X/S$ a une section] les \unsure{hom} sont définis en tous cas, mais l'exactitude demande que $\underline H^0_{\unsure{X/S|f/S}}(\underline O_X)=\underline O_S$}

l'\inserted{\unsure{dernière|dernier}}image étant formée des sections de $\underline H^1_{X/S}(\underline O_X^*)$ qui sont dans l'image de $\underline H^1(\underline O_X^*)$ i.e. dont l'obstruction dans $H^2(S,\underline O_S^*)$ est nulle

[N.B. le début de la suite exacte n'est que la suite exacte de \unsure{bas degré} de la suite spectrale de Leray

\diagram{\begin{tikzcd}[ampersand replacement=\&,column sep=tiny,row sep=small]
0\arrow[r]\&H^1(S,\underline H^0(\underline O_{X\deleted{\unsure{*}}}^*))\arrow[r]\&H^1(X,\underline O_{X\deleted{*}}^*)\arrow[r]\&H^0(S,\underline H^1(\underline O_{X\deleted{*}}^*))\arrow[r]\&H^2(S,\underline H^0(\underline O_X^*))\arrow[d]\\
{}\&{}\&{}\&{}\&H^2(X,\underline O_{X\unsure{*}}^{\unsure{*}})
\end{tikzcd}}

compte tenu que $\underline H^0(\underline O_X^*)=\underline O_S^*$]. S'il y a une section \inserted{$s$} de $X/S$, alors on \deleted{aura}\inserted{a} \deleted{un hom inverse} un résultat plus complet

\deleted{$\unsure{s^*}$.}

\uline{Corollaire 2} Soit $\underline C$ la catégorie des faisceaux inversibles $\inserted{\underline L}$ sur $X$ munis d'un \uline{isomorphisme} $\underline O_S\longrightarrow s^*(\underline L)$. Alors \deleted{\unsure{pour}}\inserted{(i)} tout automorphisme d'un objet de $\underline C$ est trivial (ii) Tout \deleted{tout} élément de $\mathrm{Prépic}\,(X/S)$ provient d'un $\underline L\in\underline C$, déterminé à isom. \uline{unique} près.

% archive sheet number: 173
